\section{Recuperación del logaritmo físico}
Sea $P_H(Y)=\sum_{l\le H}b_l[Y^{\otimes l}]$ y $D_h=Z_h\exp(L(m)-h[m])$.
Se tiene $\|D_h^{\pm1}\|_\rho\le e^{CS}$, usando $\|L(m)\|\le C(S+S^2)/j$
y $\|h\|\|m\|\le C(1+S)S/j$. Afirmamos que, para una constante $c'>0$ fija,
\[
 \|D_h^{-1} Z-E_h P_H(T-m)\|_\rho\le C e^{-c' j}.
\]
Elegimos antes del argumento el margen exponencial fijo de la parte 04 suficientemente
grande para todos los factores de tipo $C_0^v$ y $e^{Bv}$ que aparecen abajo.
Se fija $\epsilon$ del permanente, después $\kappa$ pequeña, $a\ll\sqrt\kappa$ y por último
$\eta$ pequeña respecto de esas constantes. El $a_*$ usado para $C_0$ es fijo de antemano.
\paragraph{Soportes cortos.}
En $|U|\le aj$ la comparación exacta de menores permite cambiar $p_U$ por
$\prod b_e\,e^{L(T_U)}$ con error relativo $O(e^{-cj})$. La suma de errores absolutos
normalizada cuesta $\le e^{-cj+CS}$: en este tramo $|L(T_U)|\le C(1+|U|)$, cubierto
por el peso exponencial fijo. Vale la factorización exacta
\[
 e^{L(T_U)}=e^{L(m)-h[m]}e^{h[T_U]}B(T_U-m).
\]
La norma de $Y=T_U-m$ es $\le3aj+CS$. En radio $R=c_1\sqrt\kappa j$ vale
$\sup\|B\|\le e^{C c_1^2\kappa j}$. Tomando $a/\sqrt\kappa$ y $\eta/\sqrt\kappa$
pequeños, el resto al orden $H$ es $O(e^{-c_2\kappa j})$: aplicar Cauchy
en la recta de $Y$ y sumar $(\|Y\|/R)^{H+1}e^{C R^2/j}/(1-\|Y\|/R)$.
Los factores debidos a la suma de pesos y a la división son $\le e^{CS}$,
que se absorbe reduciendo $\eta$ otra vez.
\paragraph{Soportes largos.}
Para $|U|>aj$, la cota global de $p_U$ y el peso adicional $e^{|U|}$ dan
una contribución física normalizada $\le e^{-aj+CS}$. Para el polinomio se tiene
$l!\|b_l\|\le(C\sqrt{\kappa})^l$ en $l\le H$, interpretando los grados cero y uno
por separado. Luego $\|P_H(Y)\|\le e^{C\sqrt\kappa\|Y\|}$ para cualquier $Y$,
también fuera de la bola de $L$. Se absorbe este factor por la misma reserva exponencial
usando $\|T_U-m\|\le3|U|+CS$. Queda justificada la afirmación para soportes de todo tamaño.
\paragraph{Paso al logaritmo.}
La parte 05 da $\|D_h^{-1}Z-1\|\le C(1+S)/j+Ce^{-c'j}<1/2$ para $\eta$ pequeña.
La serie del logaritmo converge en norma. En cada evaluación real $D_h>0$ porque
$Z_h=\exp G(h)$ con $G(h)$ real, y el factor que queda difiere de uno en menos de $1/2$.
Así la rama es precisamente $F$ y
\[
 \|F\|_\rho\le \|G(h)\|_\rho+\|L(m)-h[m]\|_\rho+
           \|\log(D_h^{-1}Z)\|_\rho\le C(1+S).
\]
